\chapter*{Özgeçmiş}

Yusuf Kaya, 1998 yılında Çorum'da doğdu.
Lisans eğitimini 2017--2023 yılları arasında Yıldız Teknik Üniversitesi Elektronik ve Haberleşme Mühendisliği Bölümü'nde tamamladı (AGNO: 3.03).
Lisans öğrenimi süresince TÜBİTAK 2209-A Üniversite Öğrencileri Araştırma Projeleri Destekleme Programı kapsamında Radio over Fiber (RoF) tabanlı mobil haberleşme sistemi üzerine proje yürüttü.
2024 yılında Kocaeli Üniversitesi Fen Bilimleri Enstitüsü Yazılım Mühendisliği Anabilim Dalı'nda yüksek lisans eğitimine başladı (AGNO: 3.13).

Ağustos 2022'den itibaren TÜBİTAK BİLGEM'de çalışmaya başladı; Ağustos 2022--Mayıs 2023 tarihleri arasında stajyer araştırmacı olarak OpenStack altyapısı ve ağ güvenliği alanlarında deneyim kazandı.
Mayıs 2023'ten itibaren Yazılım ve Bulut Mühendisi olarak Safir Bulut ve Multi-Cloud Management projelerinde görev almaktadır.
Bu süreçte OpenStack servis geliştirme, Kubernetes tabanlı izleme altyapısı, Prometheus/Thanos/OpenSearch entegrasyonları ve Go/Python ile servis geliştirme konularında çalışmalar yürütmüştür.

YÖKDil İngilizce puanı 75'tir (2023).
CCNA sertifikasyonunu (Ağlar, Yönlendirme ve Anahtarlama, Kurumsal Ağlar) ve Architecting on AWS sertifikasını tamamlamıştır.
